% Part VII gap-recovery problems
% Provenance:
%   VII/27, VII/28, VII/36, VII/41: CANONICAL_AUTHORED
%   VII/29, VII/30: CANONICAL_AUTHORED_FROM_SOURCE
% These are intended as insertion-ready companion problems; stable CP-VII IDs
% should be assigned only when the final Part VII ordering is frozen.

\section{Riemannian curvature contractions}

\begin{problem}[Ricci and scalar curvature of \(S^2(r)\times\mathbb{R}\)]
\label{prob:vii-gap-27}
Let
\[
M=S^2(r)\times\mathbb{R}
\]
carry the product metric.  At a point choose an orthonormal frame
\(e_1,e_2,e_3\), where \(e_1,e_2\) are tangent to \(S^2(r)\) and
\(e_3\) is tangent to the \(\mathbb{R}\)-factor.

Using the sectional curvatures of a Riemannian product:

\begin{enumerate}
\item determine \(K(e_1,e_2)\), \(K(e_1,e_3)\), and \(K(e_2,e_3)\);
\item compute \(\operatorname{Ric}(e_i,e_i)\) for \(i=1,2,3\);
\item compute the scalar curvature \(S\);
\item decide whether the product metric is Einstein.
\end{enumerate}
\end{problem}

\begin{solution}
The round sphere of radius \(r\) has constant sectional curvature
\(1/r^2\).  A two-plane containing one direction from each factor of a
Riemannian product has sectional curvature zero.  Hence
\[
K(e_1,e_2)=\frac1{r^2},
\qquad
K(e_1,e_3)=K(e_2,e_3)=0.
\]

For an orthonormal frame,
\[
\operatorname{Ric}(e_i,e_i)
 =\sum_{j\ne i}K(e_i,e_j).
\]
Therefore
\[
\operatorname{Ric}(e_1,e_1)=\frac1{r^2},
\qquad
\operatorname{Ric}(e_2,e_2)=\frac1{r^2},
\qquad
\operatorname{Ric}(e_3,e_3)=0.
\]
All mixed components vanish in this product frame, so
\[
\operatorname{Ric}
=
\operatorname{diag}\!\left(\frac1{r^2},\frac1{r^2},0\right).
\]

Taking the trace gives
\[
S=\frac2{r^2}.
\]

An Einstein metric must satisfy \(\operatorname{Ric}=\lambda g\) for a
single scalar \(\lambda\).  Here the Ricci eigenvalues are
\(1/r^2,1/r^2,0\), so no such \(\lambda\) exists.  Thus
\(S^2(r)\times\mathbb{R}\) with the product metric is not Einstein.
\end{solution}


\section{Weyl curvature}

\begin{problem}[Constant sectional curvature has zero Weyl tensor]
\label{prob:vii-gap-28}
Let \((M^n,g)\), \(n\ge4\), have constant sectional curvature \(K\), so
\[
R_{ijkl}
=
K\bigl(g_{ik}g_{jl}-g_{il}g_{jk}\bigr).
\]
Use
\[
W_{ijkl}
=
R_{ijkl}
-\frac1{n-2}
\left(
g_{ik}R_{jl}-g_{il}R_{jk}
-g_{jk}R_{il}+g_{jl}R_{ik}
\right)
+\frac{S}{(n-1)(n-2)}
\left(
g_{ik}g_{jl}-g_{il}g_{jk}
\right)
\]
to prove that \(W=0\).
\end{problem}

\begin{solution}
For constant sectional curvature,
\[
R_{ij}=(n-1)K\,g_{ij},
\qquad
S=n(n-1)K.
\]
Substitution into the middle term gives
\[
-\frac{2(n-1)K}{n-2}
\bigl(g_{ik}g_{jl}-g_{il}g_{jk}\bigr),
\]
while the scalar-curvature term becomes
\[
\frac{nK}{n-2}
\bigl(g_{ik}g_{jl}-g_{il}g_{jk}\bigr).
\]
Consequently the coefficient of
\(g_{ik}g_{jl}-g_{il}g_{jk}\) in \(W\) is
\[
K-\frac{2(n-1)K}{n-2}+\frac{nK}{n-2}=0.
\]
Hence
\[
\boxed{W=0}.
\]
Thus every constant-sectional-curvature metric is conformally flat in
dimension \(n\ge4\).
\end{solution}


\section{Lorentzian geometry}

\begin{problem}[Timelike, spacelike, and null vectors in Minkowski space]
\label{prob:vii-gap-29}
On \(\mathbb{R}^{1,3}\) use the Lorentzian metric
\[
g=-dt^2+dx^2+dy^2+dz^2.
\]
Classify each vector as timelike, spacelike, or null:
\[
u=(2,1,1,0),\qquad
v=(1,2,0,0),\qquad
w=(\sqrt2,1,1,0).
\]
Then determine the equation of the null cone through the origin.
\end{problem}

\begin{solution}
For \(X=(X^0,X^1,X^2,X^3)\),
\[
g(X,X)=-(X^0)^2+(X^1)^2+(X^2)^2+(X^3)^2.
\]
Hence
\[
g(u,u)=-4+1+1=-2<0,
\]
so \(u\) is timelike;
\[
g(v,v)=-1+4=3>0,
\]
so \(v\) is spacelike; and
\[
g(w,w)=-2+1+1=0,
\]
so \(w\) is null.

A vector \((t,x,y,z)\) is null exactly when
\[
-t^2+x^2+y^2+z^2=0.
\]
Therefore the null cone is
\[
\boxed{t^2=x^2+y^2+z^2}.
\]
Unlike a positive-definite Riemannian metric, a Lorentzian metric thus
has nonzero vectors of zero squared length.
\end{solution}


\begin{problem}[Riemannian slices and null curves in a flat FLRW metric]
\label{prob:vii-gap-30}
Consider the spatially flat FLRW metric
\[
ds^2=-dt^2+a(t)^2\bigl(dx^2+dy^2+dz^2\bigr),
\qquad a(t)>0.
\]

\begin{enumerate}
\item Find the metric induced on the hypersurface \(t=t_0\).
\item Show that for a comoving worldline \(x,y,z=\mathrm{const}\), proper
      time equals the coordinate time \(t\).
\item For a radial null curve, with
      \(dr^2=dx^2+dy^2+dz^2\), derive \(dr/dt\).
\item Explain one structural difference between the Lorentzian metric
      above and its induced Riemannian spatial metric.
\end{enumerate}
\end{problem}

\begin{solution}
On the slice \(t=t_0\) we have \(dt=0\), so
\[
ds_{\mathrm{slice}}^2
=
a(t_0)^2\bigl(dx^2+dy^2+dz^2\bigr).
\]
This is positive definite.

For a comoving worldline \(dx=dy=dz=0\), hence
\[
ds^2=-dt^2.
\]
With the convention \(ds^2=-d\tau^2\) on timelike worldlines,
\[
d\tau=dt
\]
for increasing \(t\).

For a radial null curve,
\[
0=-dt^2+a(t)^2dr^2,
\]
so
\[
\boxed{\frac{dr}{dt}=\pm\frac1{a(t)}}.
\]

The spacetime metric is indefinite and admits timelike and null
directions.  The induced metric on a constant-\(t\) spatial slice is
positive definite, so every nonzero tangent vector to the slice has
positive squared length.
\end{solution}


\section{Hyperbolic three-space}

\begin{problem}[\(\mathrm{PSL}(2,\mathbb{C})\), boundary dilation, and an isometry of \(\mathbb{H}^3\)]
\label{prob:vii-gap-36}
Use the upper-half-space model
\[
\mathbb{H}^3
=
\{(z,t)\in\mathbb{C}\times\mathbb{R}_{>0}\},
\qquad
ds^2=\frac{|dz|^2+dt^2}{t^2}.
\]
For \(s\in\mathbb{R}\), consider
\[
A_s=
\begin{pmatrix}
e^{s/2}&0\\
0&e^{-s/2}
\end{pmatrix}
\in \mathrm{SL}(2,\mathbb{C}).
\]

\begin{enumerate}
\item Compute the Möbius transformation induced by \(A_s\) on the
      boundary \(\widehat{\mathbb{C}}\).
\item Show directly that
      \[
      F_s(z,t)=(e^s z,e^s t)
      \]
      is an isometry of \(\mathbb{H}^3\).
\item Show that the vertical geodesic \(\gamma(t)=(0,t)\) is preserved,
      and compute the hyperbolic distance from \((0,1)\) to \((0,e^s)\).
\end{enumerate}
\end{problem}

\begin{solution}
The Möbius action is
\[
z\longmapsto
\frac{e^{s/2}z}{e^{-s/2}}
=e^s z.
\]

Under \(F_s\),
\[
dz\longmapsto e^s dz,
\qquad
dt\longmapsto e^s dt,
\qquad
t\longmapsto e^s t.
\]
Therefore
\[
F_s^*ds^2
=
\frac{e^{2s}|dz|^2+e^{2s}dt^2}{e^{2s}t^2}
=
\frac{|dz|^2+dt^2}{t^2}.
\]
Thus \(F_s\) is an isometry.

The set \(z=0\) is invariant under \(F_s\).  Along this vertical
geodesic the line element is
\[
ds=\frac{dt}{t}.
\]
Hence
\[
d_{\mathbb H^3}\bigl((0,1),(0,e^s)\bigr)
=
\left|\int_1^{e^s}\frac{dt}{t}\right|
=
|s|.
\]
Thus this diagonal element of \(\mathrm{PSL}(2,\mathbb C)\) acts as a
hyperbolic translation along the vertical axis, while its boundary
action is the dilation \(z\mapsto e^s z\).
\end{solution}


\section{Discrete Laplacians}

\begin{problem}[The cotangent Laplacian and its quadratic energy]
\label{prob:vii-gap-41}
Let a triangulated surface have symmetric cotangent weights
\[
w_{ij}=\frac12\bigl(\cot\alpha_{ij}+\cot\beta_{ij}\bigr)
\]
on every interior edge \(ij\).  Define the stiffness-form discrete
Laplacian \(L\) by
\[
L_{ij}=-w_{ij}\quad(i\ne j),
\qquad
L_{ii}=\sum_{j\sim i}w_{ij}.
\]

\begin{enumerate}
\item Show that every row of \(L\) sums to zero.
\item Deduce that constant vertex functions lie in the kernel of \(L\).
\item Prove
      \[
      f^{\mathsf T}Lf
      =
      \sum_{\{i,j\}}w_{ij}(f_i-f_j)^2,
      \]
      where each unoriented edge is counted once.
\item If all relevant triangles are equilateral, compute \(w_{ij}\).
\end{enumerate}
\end{problem}

\begin{solution}
For a fixed vertex \(i\),
\[
\sum_j L_{ij}
=
L_{ii}+\sum_{j\sim i}L_{ij}
=
\sum_{j\sim i}w_{ij}-\sum_{j\sim i}w_{ij}
=0.
\]
Hence, for the constant vector \(\mathbf 1\),
\[
L\mathbf 1=0.
\]

Using the symmetry \(w_{ij}=w_{ji}\),
\[
\begin{aligned}
f^{\mathsf T}Lf
&=
\sum_i\left(\sum_{j\sim i}w_{ij}\right)f_i^2
-\sum_i\sum_{j\sim i}w_{ij}f_if_j\\
&=
\sum_{\{i,j\}}w_{ij}
\bigl(f_i^2+f_j^2-2f_if_j\bigr)\\
&=
\sum_{\{i,j\}}w_{ij}(f_i-f_j)^2.
\end{aligned}
\]
Thus the stiffness energy is nonnegative whenever the cotangent weights
are nonnegative.

For equilateral triangles,
\[
\alpha_{ij}=\beta_{ij}=\frac{\pi}{3},
\qquad
\cot\frac{\pi}{3}=\frac1{\sqrt3}.
\]
Therefore
\[
\boxed{w_{ij}=\frac1{\sqrt3}}.
\]
\end{solution}
